\chapter{表象理论}

\textbf{3-1}\quad 在 $p$ 表象求解 $\delta$ 势阱 $V(x) = -\gamma \delta (x)$ 中的定态能量和波函数 $(\gamma > 0)$ .

\vfill

\newpage

\textbf{3-2}\quad 已知在 $\delta$ 势阱 $V(x) = -\gamma \delta (x)$ 中的定态归一化波函数（ $p$ 表象）为 $\varphi(p) = \frac{A}{p^2 + \hbar^2 k^2}, A = \sqrt{\frac{2\hbar^3 k^3}{\pi}}, k = \frac{\mu\gamma}{\hbar^2}$ ，试计算 $\Delta x \Delta p$ ，验证测不准关系。

\vfill

\newpage

\textbf{3-3}\quad 在 p 表象计算一维谐振子定态能量和波函数.

\vfill

\newpage

\textbf{3-4}\quad 质量为 $\mu$ 的粒子在均匀力场 $f(x) = -F(F > 0)$ 中运动, 运动的范围限制在 $x > 0$ . 给出动量表象中的定态方程, 并求出定态波函数 $\varphi(p)$ .

\vfill

\newpage

\textbf{3-5}\quad 质量为 $\mu$ 的粒子在均匀力场 $f(x) = -F(F > 0)$ 中运动, $\rho(p, t)$ 为其在动量空间的概率密度, 求 $\partial \rho / \partial t$ 与 $\partial \rho / \partial p$ 的关系.

\vfill

\newpage

\textbf{3-6}\quad 已知 $t = 0$ 时一维自由粒子波函数在坐标表象和动量表象分别是
$$
\psi (x) = N x \exp \left(- a x ^ {2} + \mathrm{i} \frac {p _ {0} x}{\hbar}\right)
$$
$$
\varphi (p) = C \left(p - p _ {0}\right) \exp \left[ - b \left(p - p _ {0}\right) ^ {2} \right]
$$
其中 $a, b$ 和 $p_0$ 都是已知实数， $N$ 与 $C$ 是归一化常数。试求 $t = 0$ 时和 $t > 0$ 时粒子坐标和动量的平均值 $\overline{x}$ 和 $\overline{p}$ 。

\vfill

\newpage

\textbf{3-7}\quad 中子 $n$ 和反中子 $\overline{n}$ 的质量都是 $m$ , 它们的态 $|n\rangle$ 和 $|\overline{n}\rangle$ 可以看成是一自由哈密顿量 $\hat{H}_0$ 的简并态: $\hat{H}_0 |n\rangle = mc^2 |n\rangle, \hat{H}_0 |\overline{n}\rangle = mc^2 |\overline{n}\rangle$ . 设有某种相互作用 $\hat{H}'$ 能使中子与反中子相互转变: $\hat{H}' |n\rangle = \alpha |\overline{n}\rangle, \hat{H}' |\overline{n}\rangle = \alpha^* |n\rangle$ , 其中 $\alpha = \alpha^*$ . 试求 $t = 0$ 时刻的一个中子在 $t$ 时刻变成反中子的概率.

\vfill

\newpage

